% =============================================================================
% 022 / 068
% Status: raw
% =============================================================================
% Notes:
%
% Source question:
% 我是之前那個問您老有無興趣和精力做影片的那位，其實我指的“比較系統的文法和發音影片”是“因為一般吳語母語者，沒有專業研究背景，最多隻熟悉自己的片區。但新吳越知識分子應該瞭解別的片區，而我本人確實也對此有很濃厚興趣。所以您老如果能出這種影片系列，最好不過，當然確實工程量較大，只是一個冒昧的詢問”。
% =============================================================================

\chapter[我是之前那個問您老有無興趣和精力做影片的那位，其實我指的“比較系統的文法和發音影片”是“因為一般吳語母語者，沒有專業研究背景，最多隻熟悉自己的片區。但新吳越知識分子應該瞭解別的片區，而我本人確實也對此有很濃厚興趣。所以您老如果能出這種影片系列，最好不過，當然確實工程量較大，只是一個冒昧的詢問”。]{我是之前那個問您老有無興趣和精力做影片的那位，其實我指的“比較系統的文法和發音影片”是“因為一般吳語母語者，沒有專業研究背景，最多隻熟悉自己的片區。但新吳越知識分子應該瞭解別的片區，而我本人確實也對此有很濃厚興趣。所以您老如果能出這種影片系列，最好不過，當然確實工程量較大，只是一個冒昧的詢問”。}
\qaanswer{
影片承載的資訊量不如一篇條理清晰的論文。學習任何東西都需要支付成本。對一門還未形成規範和標準的語言做影片教學，工程量大，效果差，事倍功半。我能夠給出的論文定期會在我的Medium上更新。學不學在於吳語區人士自己的選擇。發音這種東西還沒有形成標準，吳語拼音這種音位拼音就是一個音位對應多個音值，這樣記錄時用同樣的拼寫方式吳語區各地的人按自己的習慣發自己地方的音即可。工具網站可以參考 http://wu-chinese.com/bbs/portal.php

另外據我現在所掌握的資訊，吳語區內即使是對吳語感興趣的吳語使用者多數也還是藍青吳語使用者。他們基本上都中了「漢字中心思維」和「漢語普通話思維」很深的「毒」，多數人習慣用「漢語普通話思維」來理解吳語建構，這個毒不排乾淨，學習以原生吳語建構的吳語播音體會造成學習障礙。學吳語播音體前最好把關於漢語語境下的語言學概念和學習習慣都去除掉。首先就從「正字」、「音韻」這兩點去除。吳語的漢字根據不同時代形成的詞彙有多種讀音，而多個讀音的詞往往用到的是同一個漢字，不需要弄出一些五花八門的生僻漢字。另外要區分實詞和虛詞，虛詞不要用漢字表述。
}

